\subsection{Feedback creates an execution dependency}
Consider a service that receives an immutable audio stream and repeatedly
uses generated output as input to later recognition. A protected selection
path produces \emph{Canonical} facts. A faster \emph{Private} path may
compute and display provisional hypotheses while selection is pending.
The selected correction can change the state consumed by the next call:
StreamCorrect, for example, feeds its chosen correction back into subsequent
ASR~\cite{vu2026streamcorrect}. The problem is to coordinate this dependency
while the two paths advance asynchronously.

For a fixed session and execution contract, write an installed fact as
\begin{equation}
 C=(v,F,Z_C,P_C,T_C).
 \label{eq:canonical}
\end{equation}
Here $v$ is its version, $F$ its certified audio frontier, $Z_C$ the complete
selected output, $P_C$ the feedback state for later calls, and $T_C$ its bound
display text. A contract specifies the model, adapter, decoding configuration,
audio identity, and transformations from a complete output to $P$ and $T$.
Here $F$ is a state-processing boundary. Consumer coverage follows the
adapter's publication policy; advancing $F$ need not add visible characters.
Certification means that the protected selection and installation protocol
accepted this fact. Recognition accuracy against human transcription is a
separate property.

The adapter declares a conservative representation of \emph{all mutable
cross-call state that can affect later model inputs, selection, or feedback}.
Telemetry and device resources are recorded separately. The Qwen adapter passes prompt
tokens between calls and carries no numerical KV cache across them. The
native correction adapter additionally carries audio buffers, token groups,
context, alignment position, and decoder configuration. A state comparison
is meaningful only within its declared contract: text equality cannot stand
in for an unrecorded alignment cursor or numerical cache. The representation
may include more fields than a particular next call needs; we do not require
it to be minimal.

\subsection{Three lifetimes of provisional work}
A complete provisional node $n$ records a frozen call
\begin{equation}
 q_n=(a_n,P_n^-,b_n,\chi_n),\quad
 n=(q_n,Z_n,P_n^+,T_n),
 \label{eq:node}
\end{equation}
where $a_n$ is the input-state frontier, $b_n>a_n$ the target audio frontier,
and $\chi_n$ binds the remaining inputs and execution identity. The model
may read the complete audio prefix through $b_n$; $a_n$ does not claim a
token-to-audio alignment. $P_n^+$ and $T_n$ are derived from the complete
result through the contract's adapter. An in-flight call has a frozen input
and no complete output yet. Partially generated or cancelled tokens cannot
become a complete node.

Such work has three distinct lifetimes. Its \emph{dependency lifetime}
determines whether it can still be used under the current Canonical root.
Its \emph{visible lifetime} determines whether the Consumer keeps its whole
published hypothesis. Its \emph{physical lifetime} ends when the executor
has actually finished or drained the work. A new root can change the first
two without immediately ending the third. RevTrace connects their decisions
while keeping the evidence for each decision explicit.

The service retains every accepted input obligation until certified
processing covers it. Discarding a branch, requesting cancellation, or
recovering an old view does not remove that obligation or reset its arrival
clock. In exact mode, a missed deadline never authorizes an uncertified
Canonical transition.

\subsection{Trust and failure boundary}
We assume one trusted Owner per session, a protected selection path,
immutable accepted audio and model configuration, unambiguous encoding,
collision-resistant digests, and storage that honors successful
synchronization. Calls used in the dependency argument must actually consume
their frozen inputs and return the corresponding complete results. This
execution-integrity premise is explicit: host receipts and text hashes bind
API objects; they are not GPU arithmetic attestation. Private and Authority
need not return identical results from independent floating-point executions.

At its declared interface, the runtime checks result identity, completeness,
derived-state consistency, and ancestry, including late or repeated messages.
Provisional content can still be inaccurate. The separately frozen recovery configuration
handles whole-process termination on the same intact filesystem.
A compromised coordinator or Authority,
Byzantine host code, arbitrary persistent corruption, and lost storage are
outside this boundary. The controlled wrong-input injections in
Section~\ref{sec:eval} intentionally violate normal Private input fidelity;
we use them to study containment and executor release, not as examples
covered by the lawful-continuation proposition.
